\section{Proposed Research}
\done\todo{Introduce proposed research: (a) data curation from NVD; (b) determine vulnerability type; (c) determine prerequisite knowledge; (d) score difficult levels for learning}
\begin{figure*}[!htbp]
  \centering
  \includegraphics[width =0.9\linewidth]{figures/fig-stages.pdf}
  \caption{Four key stages in our methodology for curating NVD CVEs.}
  \label{fig:stages}
\end{figure*}

Despite holding thousands of CVEs, many with real patches, the NVD lacks in two primary areas: complete CWE labelling and an ability to sort by educational complexity. As such, instructors and student are unable to easily find examples that are suitable to their curriculum and contain correct patches. We propose a system with four key stages (Figure \ref{fig:stages}): retrieve and filter CVEs, classify by CWE, annotate prerequisite concepts, and assign difficulty levels. Our proposed research aims to transform raw NVD data into a curated set of vulnerability examples for educational use. First, we will curate and filter CVE records in order to obtain those that have complete and updated patches.  We will classify CVEs by their vulnerability types, with a specific focus on software development, as previously noted. Then, we will analyze examples for any prerequisite knowledge vital to understanding the vulnerability, which will be added in annotations. Lastly, we will determine and assign difficulty levels, such as "Beginner, Intermediate, Advanced, Expert", so that instructors and students alike can find examples suitable for any given skill level.
\subsection{Data Curation from the NVD}
We will filter for open-source projects — closed source projects do not contain patches. Then, we will write a program that will work with NVD's REST API to retrieve CVE entries published within a time window to be determined. There are several attributes of interest that we will extract from each CVE's JSON record, such as CVE ID, description, references, and timestamps — entries that lack reference URLS to patches will be dropped. Furthermore, each patch URL will be checked to ensure its functionality, and will be retrieved if they are functional. Meta data from these patches will also be extracted: at least a repository URL, revision IDs for before and after a patch, and the necessary NVD data. The result will be a curated dataset of CVE examples that contain functional patches.
\subsection{Determining Vulnerability Type}
In order to assist instructors to utilize the CVEs in this set appropriately according to their curriculum, we will classify the curated CVEs by their type. First, we will extract CWE labels present in the NVD, as well as note the CVEs that lack CWE labels entirely. For those that lack CWE labels, we will reference the CWE graph and assign the most specific label that matches the vulnerability. We will then reference the weaknesses shown in MITRE's Software Development View (ID 699), which highlights the CWEs that are most likely to occur in commonly encountered programming errors \cite{cweswview}. The view entry is not to be used as a vulnerability type -- it is meant for navigational purposes only. As such, each CVE's CWE code (e.g., CWE-119, CWE-79, CWE-489) will be translated into a corresponding readable label — such as "Buffer Overflow", "Cross-Site Scripting", or "Active Debug Code" — this will be done via a mapping tool \cite{Chen_CS_Topic-CWE-CVE_Mapping_2025}. This process will ensure that every example has a clear vulnerability type.
\subsection{Determining Prerequisite Knowledge.}
While the automation of determining prerequisite concepts remains an open question, we intend to implement a practical approach to this. Once examples receive their respective CWE label, we will first map  CWEs and their respective prerequisites using an approved mapping tool \cite{Chen_CS_Topic-CWE-CVE_Mapping_2025}, which will align examples with learning goals defined in the most recent Computer Science Curricula guidelines \cite{cscurriculum2023}. A possible next step beyond this could be the implementation of machine learning to further reduce gaps in our prerequisite listing. By using existing patches and mappings, we may be able to train a classifier, and test its accuracy for determining prerequisite knowledge. If there exists any gaps in prerequisite tagging after mapping and even possibly machine learning implementation, we intend to resort to manual tagging.
\subsection{Scoring Difficulty Levels for Learning}
As we move into the final stage — quantifying difficulty levels for learning — we acknowledge the research problems and unknowns that are possible. As such, any heuristics and machine learning methods used will undergo iterative tuning as needed. We will assign every example a didactic difficulty level (e.g., Beginner, Intermediate, Advanced, Expert), through analysis of their respective patches. We intend to measure difficulty through quantifiable characteristics of patches, such as length and amount of files impacted. We will also consider factors such as programming language, topic and the nature or platform of the software itself. These will be used to determine an overall difficulty score to be placed in appropriate buckets. This will allow instructors and students alike to filter through the dataset to find examples to match the course's respective curriculum.
